% Independent research contribution for the parent article.
% Required packages: amsmath,amssymb,amsthm; ordinary theorem environments.
% All labels and bibliography keys have qtr: prefix.
\section{The complete quartic Tornheim layer and the formal jets at every order}
\label{qtr:section}

The incoming \emph{Independent Orders and Exact Reductions} report
reduces every cubic directional derivative of the Tornheim function to
two scalar coordinates and asks for its quartic layer
\cite{qtr:Independent}. We determine that layer explicitly. Three
coordinates suffice: the diagonal fourth derivative, one convergent
logarithmic Gamma moment, and one absolutely convergent integral of
ordinary polylogarithmic order derivatives. We also determine the exact
dimension of the homogeneous jet space allowed by the source's symmetry,
axis, and Euler-plane identities. These are statements about a
specified formal relation space, not about arithmetic independence of
the constants.

\subsection{The analytic normalization and its exact restrictions}

Write
\[
 T(A,B,C)=\sum_{m,n\ge1}m^{-A}n^{-B}(m+n)^{-C},\qquad
 L=\log(2\pi),\quad z_j=\zeta^{(j)}(0),\quad
 y_j=\zeta^{(j)}(-1).
\]
The series is initially used only in its domain of absolute
convergence; elsewhere \(T\) denotes its meromorphic continuation.
The origin is not a joint regular point. For a direction
\((a,b,c)\) satisfying \((a+c)(b+c)\ne0\), put
\[
 W_{a,b,c}(t)=T(at,bt,ct),
\]
where the restriction is formed before continuation to \(t=0\).

The source constructs the unique symmetric holomorphic germ
\(R(A,B,C)=R(B,A,C)\) such that
\begin{equation}\label{qtr:polar}
 T(A,B,C)=\zeta(A)\zeta(B)
 -\frac{C\zeta(B-1)}{A+C}
 -\frac{C\zeta(A-1)}{B+C}+C R(A,B,C).
\end{equation}
For completeness, Jonqui\`ere's expansion \cite[25.12.12]{DLMF}
gives this by subtracting the following six local terms from the
Mellin kernel:
\begin{align}
 \mathcal S(A,B;x)={}&
 \Gamma(1-A)\Gamma(1-B)x^{A+B-2}
 +\Gamma(1-A)\zeta(B)x^{A-1}
 +\Gamma(1-B)\zeta(A)x^{B-1}\notag\\
 &-\Gamma(1-A)\zeta(B-1)x^A
 -\Gamma(1-B)\zeta(A-1)x^B+\zeta(A)\zeta(B).
 \label{qtr:S}
\end{align}
The subtracted remainder and every fixed parameter derivative are
normally integrable after multiplication by \(x^{C-1}\) in a
sufficiently small common polydisc. Before removing the Gamma
factors the polar terms contain
\(\Gamma(1-A)/\Gamma(1+C)\). Their difference from \(1\) is
divisible by \(A+C\), since it vanishes when \(C=-A\).
This proves that the difference between the resulting polar model
and \eqref{qtr:polar}, divided by \(C\), is holomorphic.

We will use precisely two functional restrictions of this germ:
\begin{align}
 C R(0,0,C)&=\zeta(C-1)-\zeta(C)-\frac5{12},
 \label{qtr:axis}\\
 C R(A,0,C)+A R(C,0,A)
 &=\zeta(A)\zeta(C)-\zeta(A+C)
 +\frac{\zeta(A)+\zeta(C)}2\notag\\
 &\quad+\zeta(-1)+\zeta(A-1)+\zeta(C-1).
 \label{qtr:Euler}
\end{align}
The first follows by counting pairs with prescribed sum.
The second follows from the Euler stuffle formula for
\(T(A,0,C)+T(C,0,A)\). Both are holomorphic identities near the
origin. They were proved in the incoming report; their use here
does not constitute a new proof of the cubic theorem.

\subsection{All quartic directions in three explicit coordinates}

All partial derivatives of \(R\) below are taken at the origin.
Set
\[
 \Lambda=W_{1,1,1}^{(4)}(0),\qquad
 U=R_{AAA},\qquad \Xi=R_{AAB}.
\]
The notation \(\Xi\) distinguishes this new coordinate from the
several unrelated uses of \(\kappa\) in the source articles.

\begin{lemma}[The degree-three remainder constraints]
\label{qtr:constraints}
Write
\begin{align*}
 u&=R_{AAA}=R_{BBB},& v&=R_{AAB}=R_{ABB},&
 w&=R_{AAC}=R_{BBC},\\
 x&=R_{ABC},& y&=R_{ACC}=R_{BCC},& z&=R_{CCC}.
\end{align*}
Then
\begin{align}
 z&=\frac{y_4-z_4}{4},&
 w&=\frac{z_2^2-z_4}{4},&
 u+3y&=-\frac L2z_3-z_4,\label{qtr:jetconstraints}\\
 v+x&=\frac{\Lambda+8Lz_3-12z_2^2+16z_4}{24}.
 \label{qtr:diagonalconstraint}
\end{align}
\end{lemma}

\begin{proof}
The coefficient of \(C^4\) in \eqref{qtr:axis} is \(z/6\),
and its right-hand side is \((y_4-z_4)/24\). In
\eqref{qtr:Euler}, the coefficients of \(A^2C^2\) and
\(A^3C\) are respectively \(w\) and \(u/6+y/2\).
The corresponding right-hand coefficients are
\((z_2^2-z_4)/4\) and
\((z_1z_3-z_4)/6\), where \(z_1=-L/2\).
This proves \eqref{qtr:jetconstraints}.

On the diagonal, \eqref{qtr:polar} becomes
\[
 T(t,t,t)=\zeta(t)^2-\zeta(t-1)+tR(t,t,t).
\]
Its fourth derivative is
\[
 \Lambda=-z_4-4Lz_3+6z_2^2-y_4
             +8u+24v+24w+24x+24y+4z.
\]
Substituting \eqref{qtr:jetconstraints} leaves
\(\Lambda=-8Lz_3+12z_2^2-16z_4+24(v+x)\),
as required.
\end{proof}

\begin{theorem}[Every quartic Tornheim ray]\label{qtr:quartic}
For \((a+c)(b+c)\ne0\), the restricted function \(W_{a,b,c}\)
is holomorphic at zero and
\begin{align}
 W_{a,b,c}^{(4)}(0)={}&
 abc^2\Lambda
 +4c\bigl(a^3+b^3-c^2(a+b)\bigr)U
 +12abc(a+b-2c)\Xi\notag\\
 &-2\bigl(ab(a^2+b^2-4c^2)+c^3(a+b)\bigr)Lz_3\notag\\
 &+3\bigl(2a^2b^2+c^2(a^2+b^2-4ab)\bigr)z_2^2\notag\\
 &+\left(-\frac{a^4+b^4}{2}+16abc^2
        -3c^2(a^2+b^2)-4c^3(a+b)-c^4\right)z_4\notag\\
 &+\left(c^4-\frac{cb^4}{a+c}-\frac{ca^4}{b+c}\right)y_4.
 \label{qtr:quarticformula}
\end{align}
The scalar \(U\) is given by the convergent Gamma series
\eqref{qtr:U}; \(\Xi\) is given by the convergent integral
\eqref{qtr:Xi}. No unevaluated partial derivative of \(T\)
appears in these two coordinates.
\end{theorem}

\begin{proof}
On an admissible ray the two rational polar terms in
\eqref{qtr:polar} become constant multiples of holomorphic
one-variable zeta functions, proving the asserted holomorphy.
The degree-three homogeneous part of \(R\) is
\[
 \frac{u}{6}(A^3+B^3)+\frac{v}{2}AB(A+B)
 +\frac{w}{2}C(A^2+B^2)+xABC
 +\frac{y}{2}C^2(A+B)+\frac{z}{6}C^3.
\]
Therefore the fourth derivative of \eqref{qtr:polar} is
\begin{align*}
 &-\frac{a^4+b^4}{2}z_4
 -2ab(a^2+b^2)Lz_3+6a^2b^2z_2^2
 -c\left(\frac{b^4}{a+c}+\frac{a^4}{b+c}\right)y_4\\
 &\qquad+4c\left(u(a^3+b^3)+3vab(a+b)
 +3wc(a^2+b^2)+6xabc+3yc^2(a+b)+zc^3\right).
\end{align*}
Put \(u=U,v=\Xi\) and substitute
\eqref{qtr:jetconstraints}--\eqref{qtr:diagonalconstraint}.
Collecting the five ordinary zeta combinations and the three
retained coordinates gives \eqref{qtr:quarticformula}.
\end{proof}

\paragraph{A quartic cyclic identity with no additional coordinates.}
Let \(a+b,a+c,b+c\) all be nonzero and set
\(e_1=a+b+c\), \(e_2=ab+bc+ca\), \(e_3=abc\).
Then
\begin{align}
 \sum_{\mathrm{cyc}}W_{a,b,c}^{(4)}(0)
 ={}&e_1e_3\Lambda
 +(-4e_1^2e_2+8e_2^2+12e_1e_3)Lz_3\notag\\
 &+(12e_2^2-36e_1e_3)z_2^2\notag\\
 &+(-2e_1^4+4e_1^2e_2-2e_2^2+24e_1e_3)z_4.
 \label{qtr:cyclicquartic}
\end{align}
Indeed, the cyclic coefficients of \(U\) and \(\Xi\) in
\eqref{qtr:quarticformula} vanish identically. Pairing the two
terms with each denominator cancels every \(y_4\) term.
The remaining polynomial coefficients give the displayed formula.
For example,
\begin{equation}\label{qtr:example}
 W_{1,2,3}^{(4)}(0)+W_{2,3,1}^{(4)}(0)
 +W_{3,1,2}^{(4)}(0)-36\Lambda
 =-184Lz_3+156z_2^2-386z_4.
\end{equation}
The Euler diagonal supplies an independent elementary check:
\[
 W_{1,0,1}^{(4)}(0)
 =-2Lz_3+3z_2^2-\frac{17}{2}z_4,
\]
which follows directly from
\(T(t,0,t)=(\zeta(t)^2-\zeta(2t))/2\).

More generally, the exact coordinate map
\[
 (a,b,c)\longmapsto
 \bigl(abc^2,\;
 4c(a^3+b^3-c^2(a+b)),\;
 12abc(a+b-2c)\bigr)
\]
gives a sufficient algebraic criterion for finite linear
combinations of quartic rays to reduce entirely to the ordinary
zeta derivatives displayed in \eqref{qtr:quarticformula}.
This criterion is exact in the stated three-coordinate normal
form. Additional arithmetic relations among the coordinates
could yield further reductions.

\subsection{Convergent representatives for the two additional coordinates}

Define the Stirling remainder
\[
 \rho(n)=\log\Gamma(n+1)-\left(n+\frac12\right)\log n+n
                 -\frac L2-\frac1{12n}.
\]
The source's boundary-germ theorem states, for every \(j\ge0\),
\[
 \partial_A^jR(0,0,0)
 =(-1)^j\sum_{n\ge1}(\log n)^j\rho(n)
 -y_{j+1}-\frac12z_{j+1}-y_j
 +\frac{(-1)^j\gamma_j}{12}.
\]
Consequently the coordinate needed here is
\begin{equation}\label{qtr:U}
 U=-\sum_{n\ge1}(\log n)^3\rho(n)
                  -y_4-\frac12z_4-y_3-\frac{\gamma_3}{12}.
\end{equation}
This is a specialization of an existing result, not a separate
novel Gamma-moment identity. Absolute convergence follows from
\(\rho(n)=O(n^{-3})\).

For the genuinely mixed coordinate let
\[
 f_j(x)=\left.\partial_s^j\operatorname{Li}_s(e^{-x})\right|_{s=0},
 \qquad h(x)=\log x+\gamma,\qquad j=1,2.
\]
For fixed \(x>0\), the defining series are absolutely convergent:
\[
 f_1(x)=-\sum_{n\ge1}(\log n)e^{-nx},\qquad
 f_2(x)=\sum_{n\ge1}(\log n)^2e^{-nx}.
\]
Put
\begin{align}
 P(x)={}&\frac{h(x)^3+\zeta(2)h(x)}{x^2}
 +\frac{z_1(h(x)^2+\zeta(2))+z_2h(x)}{x}\notag\\
 &-y_1(h(x)^2+\zeta(2))-y_2h(x)+z_1z_2,
 \label{qtr:P}
\end{align}
and define the ordinary absolutely convergent integral
\begin{equation}\label{qtr:I}
 I=\int_0^1\frac{f_2(x)f_1(x)-P(x)}{x}\,dx
                  +\int_1^\infty\frac{f_2(x)f_1(x)}{x}\,dx.
\end{equation}

\begin{proposition}[Explicit mixed polylogarithmic moment]
\label{qtr:Xiprop}
Let \(g=\gamma\), \(h_2=\zeta(2)\), \(h_3=\zeta(3)\), and set
\begin{align}
 E={}&-\frac{g^3+gh_2}{2}
       -\frac{3g^2+h_2}{4}-\frac{3g}{4}-\frac38\notag\\
 &-(g^2+h_2+2g+2)z_1-(g+1)z_2+gz_1z_2\notag\\
 &-\frac{g^3+3gh_2+2h_3}{3}y_1
       -\frac{g^2+h_2}{2}y_2.
 \label{qtr:E}
\end{align}
Then
\begin{equation}\label{qtr:Xi}
 \Xi=I+E.
\end{equation}
\end{proposition}

\begin{proof}
Define the normally convergent subtracted Mellin integral
\[
 I(A,B)=\int_0^1
 \frac{\operatorname{Li}_A(e^{-x})\operatorname{Li}_B(e^{-x})
                    -\mathcal S(A,B;x)}{x}\,dx
 +\int_1^\infty
 \frac{\operatorname{Li}_A(e^{-x})\operatorname{Li}_B(e^{-x})}{x}\,dx.
\]
The \(C=0\) specialization of the proof of \eqref{qtr:polar} gives
the holomorphic identity
\begin{align}
 R(A,B,0)={}&I(A,B)
 +\frac{\Gamma(1-A)\Gamma(1-B)}{A+B-2}
 +\frac{\Gamma(1-A)\zeta(B)}{A-1}
 +\frac{\Gamma(1-B)\zeta(A)}{B-1}\notag\\
 &+\gamma\zeta(A)\zeta(B)
 -\frac{\Gamma(1-A)-1}{A}\zeta(B-1)
 -\frac{\Gamma(1-B)-1}{B}\zeta(A-1).
 \label{qtr:boundarymixed}
\end{align}
Both apparent quotients at \(A=0\) or \(B=0\) are removable.
Differentiating \(\mathcal S\) twice in \(A\) and once in \(B\)
at zero produces precisely \(P(x)\). Thus
\(I_{AAB}(0,0)=I\).

To see convergence directly, write the order-differentiated
Jonqui\`ere expansion as
\[
 f_j(x)=\frac{g_j(x)}{x}
       +\sum_{k\ge0}\frac{(-1)^k\zeta^{(j)}(-k)}{k!}x^k,
 \qquad
 g_1=h,\quad g_2=h^2+\zeta(2).
\]
The subtraction \(P\) removes the singular product, both
cross terms with \(k=0,1\), and the constant analytic product.
The remaining numerator is
\(O(x(1+|\log x|^2))\) at zero. At infinity
\(f_1f_2=O(e^{-4x})\). This proves ordinary absolute convergence;
normal convergence in a parameter polydisc justifies the
differentiation used above.

Finally,
\[
 \Gamma(1-t)=1+gt+\frac{g^2+h_2}{2}t^2
               +\frac{g^3+3gh_2+2h_3}{6}t^3+O(t^4).
\]
Applying \(\partial_A^2\partial_B\) to the six elementary
terms after \(I(A,B)\) in \eqref{qtr:boundarymixed}, using this
Taylor expansion, gives exactly \(E\). This proves
\eqref{qtr:Xi}.
\end{proof}

The representation \eqref{qtr:Xi} fixes the coordinate using only
ordinary polylogarithms, their order derivatives, Gamma constants,
and ordinary zeta derivatives. It is not an asserted reduction
of \(\Xi\) to a finite algebra of standard constants.
The splitting point \(1\) in \eqref{qtr:I} is part of the formula;
changing it requires changing the explicitly integrated
subtractions.

\subsection{The formal relation space at arbitrary degree}

Here a \emph{formal degree-\(d\) freedom} means the difference of
two homogeneous degree-\(d\) jets satisfying symmetry in \(A,B\),
the same degree of \eqref{qtr:axis}, and the same degree of
\eqref{qtr:Euler}. No other functional relation is imposed.

\begin{theorem}[Exact homogeneous kernel]\label{qtr:kernel}
For \(d\ge2\), every formal degree-\(d\) freedom has a unique
representation
\begin{equation}\label{qtr:kernelform}
 \Delta R_d(A,B,C)=
 A(A-C)S(A,C)+B(B-C)S(B,C)+ABQ(A,B,C),
\end{equation}
where \(S\) is a homogeneous symmetric polynomial of degree
\(d-2\) in two variables and \(Q\) is a homogeneous polynomial of
degree \(d-2\), symmetric in \(A,B\). Therefore the dimension of
this formal freedom is
\begin{equation}\label{qtr:dimension}
 \dim\mathcal K_d
 =\left\lfloor\frac d2\right\rfloor
     +\left\lfloor\frac{d^2}{4}\right\rfloor.
\end{equation}
At degrees \(0\) and \(1\) the freedom is zero.
\end{theorem}

\begin{proof}
Put \(D(A,C)=\Delta R_d(A,0,C)\). The axis condition gives
\(D(0,C)=0\), so \(D=A V(A,C)\). The homogeneous Euler
condition becomes
\[
 AC\bigl(V(A,C)+V(C,A)\bigr)=0.
\]
Hence \(V\) is antisymmetric in \(A,C\), and
\(V(A,C)=(A-C)S(A,C)\), where \(S\) is symmetric.
This proves the asserted boundary value.

Subtract the first two terms of \eqref{qtr:kernelform} from
\(\Delta R_d\). The difference is still symmetric in \(A,B\)
and vanishes when \(A=0\) or \(B=0\). It is therefore divisible
by \(AB\); its quotient is the required symmetric polynomial
\(Q\). This construction determines \(S\) from the boundary
and then determines \(Q\), proving uniqueness. Conversely, a
direct substitution shows that every expression
\eqref{qtr:kernelform} satisfies all three homogeneous
conditions.

The space of symmetric polynomials of degree \(d-2\) in two
variables has dimension \(\lfloor d/2\rfloor\). The space
of degree-\(d-2\) polynomials in three variables symmetric
in the first two has dimension
\[
 \sum_{k=0}^{d-2}
 \left(\left\lfloor\frac{d-2-k}{2}\right\rfloor+1\right)
 =\left\lfloor\frac{d^2}{4}\right\rfloor.
\]
For \(d=0,1\), the divisibility argument forces \(S\) and \(Q\)
to vanish.
\end{proof}

The Tornheim derivative of order \(n\) uses degree \(d=n-1\)
of \(R\). Thus the first formal dimensions, for derivative
orders \(n=1,\ldots,8\), are
\[
 0,\quad0,\quad2,\quad3,\quad6,\quad8,\quad12,\quad15.
\]
In particular the quartic three-coordinate count is exact under
these stated relations. This count does not say that three
arithmetically independent periods are necessary.

\begin{corollary}[All-order cyclic cancellation and its exact image]
\label{qtr:cyclicimage}
For \(d\ge2\), the image of the homogeneous freedom under
\[
 \Delta R_d\longmapsto
 c\Delta R_d(a,b,c)+a\Delta R_d(b,c,a)
                         +b\Delta R_d(c,a,b)
\]
is exactly \(abc\) times the space of symmetric homogeneous
polynomials of degree \(d-2\) in \(a,b,c\). Consequently its
dimension is
\[
 p_{\le3}(d-2)=
 \left\lfloor\frac{(d+1)^2+3}{12}\right\rfloor.
\]
Equivalently, at derivative order \(n\ge3\), cyclic
symmetrization leaves exactly
\(\lfloor(n^2+3)/12\rfloor\) formal coordinates.
\end{corollary}

\begin{proof}
The terms involving \(S\) in \eqref{qtr:kernelform} cancel
pairwise because \(S(a,c)=S(c,a)\). The remaining expression is
\[
 abc\bigl(Q(a,b,c)+Q(b,c,a)+Q(c,a,b)\bigr).
\]
The parenthesis is fully symmetric. Every symmetric
homogeneous polynomial \(P\) occurs by choosing \(Q=P/3\),
so the map is onto. A basis of the symmetric space is
\(e_1^i e_2^j e_3^k\) with \(i+2j+3k=d-2\).
Counting these triples gives the stated dimension.
\end{proof}

\paragraph{Why the diagonal ceases to determine cyclic fifth derivatives.}
The preceding dimensions equal one at derivative orders three
and four, which explains why the source's cubic cyclic formula
and \eqref{qtr:cyclicquartic} each need just the corresponding
diagonal derivative. At order five the dimension is two.
This is not merely a count: the allowed deformation
\[
 \Delta R_4(A,B,C)=\frac{AB}{3}
 \bigl((A+B+C)^2-3(AB+BC+CA)\bigr)
\]
vanishes on the diagonal and preserves the axis and Euler
relations, but changes the cyclic fifth derivative by
\[
 5!\,abc\bigl((a+b+c)^2-3(ab+bc+ca)\bigr).
\]
Therefore the diagonal fifth derivative together with the
named relations cannot determine every cyclic fifth ray.
Additional analytic information is needed.

\subsection{An exact symmetric generating germ}

The preceding cyclic statement can be expressed without
choosing coordinates for each degree.

\begin{proposition}[Cyclic factorization]\label{qtr:factorization}
There is a unique symmetric holomorphic germ \(J(A,B,C)\)
near the origin such that
\begin{align}
 \sum_{\mathrm{cyc}}T(A,B,C)
 ={}&2\bigl(\zeta(A)\zeta(B)+\zeta(B)\zeta(C)
                              +\zeta(C)\zeta(A)\bigr)\notag\\
 &-\zeta(A+B)-\zeta(B+C)-\zeta(C+A)\notag\\
 &+2\bigl(\zeta(A)+\zeta(B)+\zeta(C)\bigr)+1
 +ABC J(A,B,C).
 \label{qtr:J}
\end{align}
\end{proposition}

\begin{proof}
In the cyclic sum of \eqref{qtr:polar}, the two terms with
denominator \(A+C\) combine to \(-\zeta(B-1)\), and likewise
for the other two denominators. Thus the cyclic sum is
holomorphic and symmetric near the origin.
On \(B=0\), the Euler relation gives
\[
 T(A,0,C)+T(0,C,A)+T(C,A,0)
                   =2\zeta(A)\zeta(C)-\zeta(A+C).
\]
The explicit zeta expression in \eqref{qtr:J} has exactly the
same restriction. By symmetry, their difference vanishes on
all three coordinate planes. The convergent Taylor series is
therefore divisible by \(ABC\). Its quotient is holomorphic
and symmetric, and uniqueness is immediate.
\end{proof}

Let \(J_k\) denote the homogeneous degree-\(k\) Taylor
polynomial of \(J\). For every \(n\ge3\) and every direction
whose pairwise sums are nonzero, differentiating the restricted
identity gives the all-order formula
\begin{align}
 \sum_{\mathrm{cyc}}W_{a,b,c}^{(n)}(0)
 ={}&2\sum_{\{u,v\}\in\{\{a,b\},\{b,c\},\{c,a\}\}}
       \sum_{j=1}^{n-1}\binom nj u^jv^{n-j}z_jz_{n-j}\notag\\
 &-\bigl((a+b)^n+(b+c)^n+(c+a)^n\bigr)z_n
 +n!\,abc\,J_{n-3}(a,b,c).
 \label{qtr:allcyclic}
\end{align}
The inner sum is symmetric in \(u,v\), so the unordered-pair
notation is unambiguous. All derivatives at \(-1\) and all
boundary Gamma coordinates have disappeared. At the first
two orders of \(J\),
\[
 J_0=\frac{\Omega}{2}+3Lz_2+4z_3,\qquad
 J_1(A,B,C)=\frac{A+B+C}{24}
                  (\Lambda+8Lz_3-12z_2^2+16z_4),
\]
where \(\Omega=W_{1,1,1}'''(0)\) is the incoming cubic
diagonal coordinate. Thus \eqref{qtr:allcyclic} recovers the
source's cubic identity and proves
\eqref{qtr:cyclicquartic}.

\subsection{Verification, scope, and the next exact questions}

The accompanying \texttt{check\_quartic\_tornheim.py} verifies
\eqref{qtr:quarticformula}, \eqref{qtr:cyclicquartic},
the Euler specialization, and the elementary correction \(E\)
by exact symbolic arithmetic. It independently constructs the
homogeneous constraint matrices for degrees \(0\) through \(10\),
computes their exact ranks and cyclic images, and checks them
against the proved dimension formulas.
An optional numerical mode evaluates the new coordinates
through integrated Jonqui\`ere series and a Stirling-subtracted
Gamma sum, and compares asymmetric fourth derivatives with
an independent Mellin evaluation of \(T\). Numerical
discrepancies are diagnostics, not interval certificates.

The results answer the quartic directional-layer question and
the formal homogeneous-dimension part of the all-order
question in \cite{qtr:Independent}. They do not evaluate
\(\Lambda\), \(\Xi\), or the earlier cubic coordinates in
a finite standard-constant algebra, and they do not settle
\(S_6\) or the revised \(S_8\).
The Mellin and Jonqui\`ere ingredients are classical
\cite{DLMF}; the quartic and formal kernel results are
developments relative to the inspected repository snapshot.
A global priority claim is not made.

The next concrete problems are to identify a useful
independent analytic relation for the symmetric degree-two
germ \(J_2\); to construct an economical Gamma or Barnes
representation for \(\Xi\); and to compare additional
Tornheim functional relations with the kernel
\eqref{qtr:kernelform}. Any improved coordinate count should
specify exactly which extra relations are used.
